\documentclass{article}%
\usepackage[T1]{fontenc}%
\usepackage[utf8]{inputenc}%
\usepackage{lmodern}%
\usepackage{textcomp}%
\usepackage{lastpage}%
%
\usepackage{amsmath}%
\usepackage{amssymb}%
\usepackage{amsfonts}%
\usepackage{graphicx}%
\usepackage{hyperref}%
\usepackage{tikz}%
\usepackage{pgfplots}%
\usetikzlibrary{math}%
\title{ITA 2027 - Matemática\thanks{This exam was solved and automatically generated with gpt-resolve: \url{https://github.com/lgabs/gpt-resolve}}}%
\author{claude{-}sonnet{-}5{-}5 (solver) + claude{-}sonnet{-}5{-}5 (judge)}%
\date{November 09, 2026}%
%
\begin{document}%
\normalsize%
\maketitle%
\section*{Questão 1}

\section*{Solução}

\textbf{Notação e condição inicial.} Seja $(a_1,a_2,\dots,a_{10})$ uma PA de razão $r$, com $a_n=a_1+(n-1)r$ e $a_n\in\mathbb{N}=\{1,2,3,\dots\}$ para todo $n$. Como $r=a_2-a_1$ é diferença de dois naturais, $r\in\mathbb{Z}$.

\textbf{Passo 1: a condição da soma.} Pela fórmula da soma dos termos de uma PA,
\[
S_{10}=\frac{(a_1+a_{10})\cdot 10}{2}=5\,(2a_1+9r)=435
\quad\Longleftrightarrow\quad 2a_1+9r=87 .
\]
Equivalentemente, $a_1+a_{10}=87$.

\textbf{Passo 2: paridade.} Como $2a_1$ é par e $87$ é ímpar, $9r$ deve ser ímpar, isto é, \emph{$r$ é ímpar}. Em particular $r\neq 0$ (a PA constante exigiria $2a_1=87$, impossível). Reciprocamente, se $r$ é ímpar, então $87-9r$ é par e $a_1=\dfrac{87-9r}{2}$ é inteiro.

\textbf{Passo 3: redução a $r>0$.} Lendo uma PA de trás para frente obtém-se outra PA, com os mesmos termos, a mesma soma e razão $-r$. Basta, então, estudar $r>0$ e depois incluir as sequências invertidas.

\textbf{Passo 4: todos os termos naturais.} Para $r>0$ a PA é crescente, e o menor termo é $a_1$. Logo a condição é $a_1\ge 1$:
\[
\frac{87-9r}{2}\ge 1\iff 9r\le 85\iff r\le \frac{85}{9}\approx 9{,}4 .
\]
Com $r$ ímpar e positivo, resta $r\in\{1,3,5,7,9\}$, e $a_1=\dfrac{87-9r}{2}$ é inteiro positivo em cada caso. (Se $0$ fosse considerado natural, a condição $a_1\ge 0$ daria o mesmo conjunto, pois $r=11$ daria $a_1<0$.)

\textbf{Passo 5: as progressões.}
\[
\begin{array}{c|c|c|c}
r & a_1 & a_{10}=a_1+9r & \text{PA}\\ \hline
1 & 39 & 48 & 39,40,41,\dots,48\\
3 & 30 & 57 & 30,33,36,\dots,57\\
5 & 21 & 66 & 21,26,31,\dots,66\\
7 & 12 & 75 & 12,19,26,\dots,75\\
9 & 3 & 84 & 3,12,21,\dots,84
\end{array}
\]

\textbf{Verificação.} Em cada linha $a_1+a_{10}=87$, logo $S_{10}=5\cdot 87=435$. Todos os termos são naturais, pois $a_1\ge 3\ge 1$ e $r>0$.

\textbf{Conclusão.} As PAs pedidas têm razão $r$ ímpar com $|r|\le 9$ e termo geral
\[
a_n=\frac{87-9r}{2}+(n-1)\,r,\qquad r\in\{1,3,5,7,9\}\ \text{(crescentes)}.
\]
Há exatamente \textbf{5 progressões crescentes}, as da tabela. Se o sentido da sequência for levado em conta, há também as \textbf{5 decrescentes} correspondentes, obtidas invertendo a ordem dos termos. Elas têm razões $-1,-3,-5,-7,-9$ e primeiros termos $48,57,66,75,84$. Não existe PA constante. Não há outras soluções.

ANSWER: As PAs são $a_n=\frac{87-9r}{2}+(n-1)r$ com $r\in\{1,3,5,7,9\}$, isto é, $(39,\dots,48)$, $(30,\dots,57)$, $(21,\dots,66)$, $(12,\dots,75)$ e $(3,\dots,84)$, com razões $1,3,5,7,9$. Cada uma pode ser lida também em ordem inversa, com razão $-r$. São 5 PAs crescentes, ou 10 contando as decrescentes.%
\newpage%
\section*{Questão 2}

\section*{Solução}

\textbf{Passo 1 (retirar o módulo do logaritmo).} Para todo $x\in\mathbb{R}$ temos $2^x>0$, logo $2^x+2>2$ e, como $\log_2$ é crescente,
\[
\log_2(2^x+2)>\log_2 2=1>0 .
\]
Assim $|\log_2(2^x+2)|=\log_2(2^x+2)$ e a desigualdade equivale a
\[
\log_2(2^x+2)>|2x| .
\]
Como $|2x|\ge 0$ e $\log_2$ é crescente, aplicar $2^{(\cdot)}$ nos dois lados preserva o sentido da desigualdade:
\[
2^x+2>2^{|2x|}.
\]
Façamos $t=2^x>0$ e separemos em casos conforme o sinal de $x$.

\textbf{Passo 2: caso $x\ge 0$.} Aqui $|2x|=2x$ e $2^{2x}=t^2$, com $t\ge 1$. A desigualdade fica
\[
t+2>t^2 \iff t^2-t-2<0 \iff (t-2)(t+1)<0 .
\]
Como $t+1>0$, isso equivale a $t<2$, isto é, $2^x<2$, ou seja, $x<1$. Combinando com $x\ge 0$:
\[
0\le x<1 .
\]

\textbf{Passo 3: caso $x<0$.} Aqui $|2x|=-2x$ e $2^{-2x}=t^{-2}$, com $0<t<1$. A desigualdade fica
\[
t+2>\frac1{t^2}\iff t^3+2t^2-1>0 \qquad (\text{multiplicamos por } t^2>0).
\]
Observando que $t=-1$ é raiz, fatoramos:
\[
t^3+2t^2-1=(t+1)(t^2+t-1).
\]
Como $t+1>0$, resta $t^2+t-1>0$. As raízes de $t^2+t-1$ são $\dfrac{-1\pm\sqrt5}{2}$; como $t>0$, a condição é
\[
t>\frac{\sqrt5-1}{2}.
\]
Note que $0<\dfrac{\sqrt5-1}{2}<1$ (pois $2<\sqrt5<3$), então esse limite está dentro do intervalo $0<t<1$. Como $t=2^x$ é crescente,
\[
x>\log_2\!\left(\frac{\sqrt5-1}{2}\right)=-\log_2\!\left(\frac{1+\sqrt5}{2}\right),
\]
onde usamos $\dfrac{\sqrt5-1}{2}=\dfrac{1}{(1+\sqrt5)/2}$ (pois $(\sqrt5-1)(\sqrt5+1)=4$). Esse número é negativo ($\approx -0{,}694$). Combinando com $x<0$:
\[
-\log_2\!\left(\frac{1+\sqrt5}{2}\right)<x<0 .
\]

\textbf{Passo 4 (união e verificação).} Unindo os dois casos (o ponto $x=0$ pertence à solução, pois $\log_2 3>0$):
\[
-\log_2\!\left(\frac{1+\sqrt5}{2}\right)<x<1 .
\]
\emph{Conferência.} Em $x=1$: $\log_2 4=2=|2x|$, igualdade, logo $x=1$ não pertence à solução (desigualdade estrita). No outro extremo, $t=\frac{\sqrt5-1}{2}$ satisfaz $t^2+t=1$, donde $t+2=t^{-2}$ e vale também a igualdade, logo o extremo é excluído. Testes: $x=-0{,}69$ dá $\log_2(2^{-0,69}+2)\approx1{,}390>1{,}38$ (vale); $x=-0{,}70$ dá $\approx1{,}387<1{,}40$ (não vale); $x=0{,}5$ dá $\log_2(2^{0,5}+2)\approx1{,}77>1$ (vale); $x=2$ dá $\log_2 6\approx2{,}58<4$ (não vale). Uma varredura numérica em $[-5,5]$ confirmou o intervalo.

Conjunto solução:
\[
S=\left(\log_2\frac{\sqrt5-1}{2},\,1\right)=\left(-\log_2\frac{1+\sqrt5}{2},\,1\right)\approx(-0{,}694;\;1).
\]

ANSWER: $-\log_2\!\left(\dfrac{1+\sqrt5}{2}\right)<x<1$, isto é, $x\in\left(\log_2\dfrac{\sqrt5-1}{2},\,1\right)$.%
\newpage%
\section*{Questão 3}

\section*{Solução}

\textbf{1) Raízes da equação.} Como $\dfrac{-1}{i}=i$, temos
\[
r+iz^8=0 \iff z^8=-\frac{r}{i}=ir=r\left(\cos\frac{\pi}{2}+i\,\mathrm{sen}\,\frac{\pi}{2}\right).
\]
Pela fórmula de De Moivre, as oito raízes são
\[
z_k=r^{1/8}\left[\cos\!\left(\frac{\pi}{16}+\frac{k\pi}{4}\right)+i\,\mathrm{sen}\!\left(\frac{\pi}{16}+\frac{k\pi}{4}\right)\right],\qquad k=0,1,\dots,7 .
\]
Todas têm o mesmo módulo $R=r^{1/8}$ e seus argumentos diferem de $\frac{2\pi}{8}=\frac{\pi}{4}$ consecutivamente. Logo, os vértices pertencem à circunferência de centro na origem e raio $R$, igualmente espaçados: o polígono é um \textbf{octógono regular} inscrito nessa circunferência.

\textbf{2) Área em função de $R$.} O octógono se decompõe em $8$ triângulos isósceles de vértice na origem, com lados $R$, $R$ e ângulo central $\frac{\pi}{4}$. Assim,
\[
A=8\cdot\frac12 R^2\,\mathrm{sen}\frac{\pi}{4}=4R^2\cdot\frac{\sqrt2}{2}=2\sqrt2\,R^2 .
\]
Com $A=\sqrt2$:
\[
2\sqrt2\,R^2=\sqrt2 \;\Longrightarrow\; R^2=\frac12 \;\Longrightarrow\; R=\frac{\sqrt2}{2}\quad(R>0).
\]

\textbf{3) Valor de $r$.} Como $R=r^{1/8}$ e $r>0$,
\[
r=R^8=\left(R^2\right)^4=\left(\frac12\right)^4=\frac1{16}.
\]

\textbf{4) Lado.} Pela lei dos cossenos no triângulo central, com ângulo $\frac{\pi}{4}$ entre dois lados de medida $R$:
\[
\ell^2=R^2+R^2-2R^2\cos\frac{\pi}{4}=2R^2\left(1-\frac{\sqrt2}{2}\right)=1-\frac{\sqrt2}{2}=\frac{2-\sqrt2}{2},
\]
pois $2R^2=1$. Portanto
\[
\ell=\sqrt{\frac{2-\sqrt2}{2}}=\frac{\sqrt{4-2\sqrt2}}{2}\approx 0{,}541 .
\]

\textbf{5) Conferência.} A área de um octógono regular de lado $\ell$ é $2(1+\sqrt2)\ell^2$. Então
\[
2(1+\sqrt2)\cdot\frac{2-\sqrt2}{2}=(1+\sqrt2)(2-\sqrt2)=2-\sqrt2+2\sqrt2-2=\sqrt2,
\]
o que coincide com o dado. Além disso, $|z_k|^8=R^8=\frac1{16}=r$, consistente com $|z^8|=r$. Valores plausíveis: $R\approx0{,}707$ e $\ell\approx 0{,}541<2R\approx1{,}414$ (diâmetro).

\textbf{Resultado:}
\[
\ell=\sqrt{\frac{2-\sqrt2}{2}}=\frac{\sqrt{4-2\sqrt2}}{2},\qquad r=\frac1{16}.
\]

ANSWER: $\ell=\sqrt{\dfrac{2-\sqrt2}{2}}=\dfrac{\sqrt{4-2\sqrt2}}{2}\approx 0{,}541$ e $r=\dfrac1{16}$%
\newpage%
\section*{Questão 4}

\section*{Solução}

\textbf{1. Fatoração de $q(x)$.} Como $q(-1)=-1+3-2=0$, dividindo por $x+1$ obtemos
\[
q(x)=x^3-3x-2=(x+1)(x^2-x-2)=(x+1)^2(x-2).
\]
Assim, $q$ tem a raiz dupla $-1$ e a raiz simples $2$.

\textbf{2. Determinação de $a$.} Se $q\mid p$, toda raiz de $q$ é raiz de $p$, com multiplicidade pelo menos igual à que tem em $q$. Logo $p(2)=0$, $p(-1)=0$ e, como $-1$ é raiz dupla, também $p'(-1)=0$.

\begin{itemize}
\item $p(2)=128+192a-32-80a+56-28a+18+a=170+85a=0\ \Rightarrow\ a=-2$.
\item $p(-1)=-1+3a+1-5a-7-7a-9+a=-16-8a=0\ \Rightarrow\ a=-2$.
\item $p'(x)=7x^6+18ax^5-5x^4-20ax^3+21x^2-14ax+9$, logo $p'(-1)=32+16a=0\ \Rightarrow\ a=-2$.
\end{itemize}

As três condições são compatíveis e dão $a=-2$. Então
\[
p(x)=x^7-6x^6-x^5+10x^4+7x^3+14x^2+9x-2 .
\]

\textbf{3. Quociente.} Dividindo $p$ por $q(x)=x^3-3x-2$ (grau do quociente $=7-3=4$):
\[
p(x)=(x^3-3x-2)\,(x^4-6x^3+2x^2-6x+1),
\]
com resto nulo. Conferindo: $x^4\cdot q=x^7-3x^5-2x^4$; $-6x^3\cdot q=-6x^6+18x^4+12x^3$; $2x^2\cdot q=2x^5-6x^3-4x^2$; $-6x\cdot q=-6x^4+18x^2+12x$; $1\cdot q=x^3-3x-2$. A soma dá exatamente $p(x)$ com $a=-2$.

\textbf{4. Raízes do quociente.} Seja $s(x)=x^4-6x^3+2x^2-6x+1$. Ele é recíproco (coeficientes palindrômicos) e $s(0)=1\neq0$. Dividindo por $x^2$ e fazendo $t=x+\dfrac1x$, de modo que $x^2+\dfrac1{x^2}=t^2-2$:
\[
\left(x^2+\frac1{x^2}\right)-6\left(x+\frac1x\right)+2=0
\iff t^2-2-6t+2=0\iff t(t-6)=0 .
\]
\begin{itemize}
\item $t=0$: $x^2+1=0\Rightarrow x=\pm i$.
\item $t=6$: $x^2-6x+1=0\Rightarrow x=3\pm2\sqrt2$.
\end{itemize}

\textbf{5. Multiplicidades.} Como $s(-1)=16\neq0$ e $s(2)=-35\neq0$, as raízes de $q$ não se repetem em $s$. Portanto $-1$ tem multiplicidade $2$ e $2$ tem multiplicidade $1$. As sete raízes contadas com multiplicidade são $-1,-1,2,i,-i,3+2\sqrt2,3-2\sqrt2$, e
\[
p(x)=(x+1)^2(x-2)(x^2+1)(x^2-6x+1).
\]

\textbf{Verificação.} Multiplicando os fatores, obtém-se $x^7-6x^6-x^5+10x^4+7x^3+14x^2+9x-2$, que é $p$ com $a=-2$. Numericamente, $p(3\pm2\sqrt2)=0$ também.

\medskip
Para $a=-2$, as raízes de $p$ são $-1$ (dupla), $2$, $i$, $-i$, $3+2\sqrt2$ e $3-2\sqrt2$.

ANSWER: $a=-2$; raízes de $p$: $-1$ (multiplicidade 2), $2$, $\pm i$ e $3\pm2\sqrt2$.%
\newpage%
\section*{Questão 5}

\section*{Solução}

\textbf{1) Vértices $A$ e $B$.} Com $x=1$ na elipse $2x^2+y^2=6$:
\[
2+y^2=6 \;\Rightarrow\; y=\pm 2 .
\]
Como $A\neq B$, podemos tomar $A=(1,2)$ e $B=(1,-2)$, de modo que $m(\overline{AB})=4$.

\textbf{2) Condição do perímetro.} Como o perímetro é $10$,
\[
m(\overline{AC})+m(\overline{BC})=10-4=6 .
\]

\textbf{3) Lugar geométrico de $C$.} Os pontos $P$ com $PA+PB=6$ (constante maior que $AB=4$) formam a elipse de focos $A$ e $B$ e eixo maior $2a'=6$. Seus parâmetros são:
\begin{itemize}
\item centro: ponto médio de $\overline{AB}$, isto é, $(1,0)$;
\item $a'=3$ e distância focal $2c'=4$, logo $c'=2$;
\item $b'^2=a'^2-c'^2=9-4=5$;
\item eixo maior vertical, pois os focos estão na reta $x=1$.
\end{itemize}
Portanto
\[
\frac{(x-1)^2}{5}+\frac{y^2}{9}=1 .
\]
Assim, $C$ pertence a esta elipse e também à elipse dada, $2x^2+y^2=6$.

\textbf{4) Interseção.} Da elipse dada, $y^2=6-2x^2$. Substituindo na equação do passo 3 e multiplicando por $45$:
\[
9(x-1)^2+5(6-2x^2)=45
\;\Longrightarrow\; 9x^2-18x+9+30-10x^2=45
\;\Longrightarrow\; x^2+18x+6=0 .
\]
Logo
\[
x=-9\pm\sqrt{81-6}=-9\pm 5\sqrt3 .
\]

\textbf{5) Escolha da raiz.} Como $C$ está na elipse $2x^2+y^2=6$, devemos ter $2x^2\le 6$, isto é, $|x|\le\sqrt3\approx1{,}73$.
\begin{itemize}
\item $x=-9-5\sqrt3\approx-17{,}7$ não serve, pois $|x|>\sqrt3$.
\item $x=5\sqrt3-9\approx-0{,}34$ serve. Nesse caso $x^2=156-90\sqrt3\approx0{,}115$ e $y^2=6-2x^2\approx5{,}77>0$, portanto existem dois pontos $C=(x,\pm y)$ com essa abscissa.
\end{itemize}
Além disso, $x\neq1$, logo $C$ não é colinear com $A$ e $B$ e o triângulo é não degenerado.

\textbf{6) Verificação numérica.} Para $x\approx-0{,}3397$ e $y\approx2{,}4019$:
\[
AC\approx\sqrt{1{,}3397^2+0{,}4019^2}\approx1{,}399,\qquad
BC\approx\sqrt{1{,}3397^2+4{,}4019^2}\approx4{,}601 .
\]
Então $AC+BC\approx6{,}000$ e o perímetro é $\approx 6+4=10$. O caso $y<0$ é simétrico e dá o mesmo resultado.

\textbf{Resposta:} a abscissa de $C$ é
\[
x_C=5\sqrt3-9\approx -0{,}34 .
\]

ANSWER: $5\sqrt{3}-9$
%
\newpage%
\section*{Questão 6}

\section*{Solução}

\textbf{Convenção.} Como $\mathbb{N}=\{1,2,3,\dots\}$, todas as variáveis são inteiras e $x_i\ge 1$.

\textbf{1. Redução do sistema.} As variáveis $x_3,x_4,x_5$ aparecem sempre na combinação $x_3+x_4+x_5$. Definimos $S=x_3+x_4+x_5$. O sistema fica
\[
\begin{cases}
x_1+x_2+S+3x_6=30 & (\mathrm{I})\\
x_1+2x_2+2S+4x_6=52 & (\mathrm{II})\\
2x_1+2x_2+3S+7x_6=76 & (\mathrm{III})
\end{cases}
\]
São 3 equações nas 4 incógnitas $x_1,x_2,S,x_6$, o que deixa um grau de liberdade.

\textbf{2. Resolução em função de $x_6$.}
\begin{itemize}
\item $(\mathrm{II})-(\mathrm{I})$: $\;x_2+S+x_6=22$, ou seja, $x_2=22-S-x_6$.
\item Substituindo em $(\mathrm{I})$: $x_1=30-x_2-S-3x_6=30-(22-S-x_6)-S-3x_6=8-2x_6$.
\item Substituindo $x_1$ e $x_2$ em $(\mathrm{III})$:
\[
2(8-2x_6)+2(22-S-x_6)+3S+7x_6=60+S+x_6=76 \;\Rightarrow\; S+x_6=16 .
\]
\end{itemize}
Logo
\[
S=16-x_6,\qquad x_2=22-16=6,\qquad x_1=8-2x_6 .
\]

\textbf{3. Restrições.} Precisamos de $x_6\ge1$ e $x_1=8-2x_6\ge 1$, isto é, $x_6\le 3$ (pois $x_6$ é inteiro). Também $x_2=6\ge 1$ vale sempre. Como $x_3,x_4,x_5\ge1$, é necessário $S\ge 3$, e $S=16-x_6\ge 13$ vale para $x_6\le 3$. Portanto $x_6\in\{1,2,3\}$.

\textbf{4. Contagem.} Fixado $x_6$, os valores $x_1$, $x_2$ e $S$ ficam determinados. Resta contar as ternas $(x_3,x_4,x_5)$ de inteiros positivos com soma $S$. Pelo método das ``estrelas e barras'', o número delas é $\binom{S-1}{2}$.

\[
\begin{array}{c|c|c|c|c}
x_6 & x_1 & x_2 & S & \binom{S-1}{2}\\\hline
1 & 6 & 6 & 15 & \binom{14}{2}=91\\
2 & 4 & 6 & 14 & \binom{13}{2}=78\\
3 & 2 & 6 & 13 & \binom{12}{2}=66
\end{array}
\]

\textbf{5. Total.}
\[
91+78+66=235 .
\]

\textbf{Verificação.} Para $x_6=1$, $x_1=x_2=6$ e $S=15$:
$6+6+15+3=30$; $6+12+30+4=52$; $12+12+45+7=76$. As três equações são satisfeitas. Uma contagem por força bruta no computador também resultou em $235$ sêxtuplas.

ANSWER: $235$%
\newpage%
\section*{Questão 7}

\section*{Solução}

Seja $D(n)=n^2-2n-27=(n-1)^2-28$. Para $n$ inteiro, $D(n)\neq 0$, pois $(n-1)^2=28$ não tem solução inteira (28 não é quadrado perfeito). Logo a fração está definida para todo $n\in\mathbb{Z}$.

\textbf{1) Divisão de polinômios.} Calculando
\[
(2n-1)(n^2-2n-27)=2n^3-4n^2-54n-n^2+2n+27=2n^3-5n^2-52n+27,
\]
obtemos
\[
2n^3-5n^2-52n+90=(2n-1)(n^2-2n-27)+63 .
\]
Portanto
\[
\frac{2n^3-5n^2-52n+90}{n^2-2n-27}=2n-1+\frac{63}{n^2-2n-27}.
\]

\textbf{2) Condição de integralidade.} Como $2n-1$ é inteiro para todo $n$ inteiro, o quociente é inteiro se, e somente se, $\dfrac{63}{D(n)}$ é inteiro, isto é, se, e somente se,
\[
D(n)\ \text{divide}\ 63,\qquad D(n)\in\{\pm1,\pm3,\pm7,\pm9,\pm21,\pm63\}.
\]

\textbf{3) Teste dos divisores.} Como $D(n)=(n-1)^2-28$, temos $(n-1)^2=28+D(n)$, que deve ser quadrado perfeito (e não negativo). Analisando cada divisor $d$ de $63$:

\begin{center}
\begin{tabular}{c|cccccccccccc}
$d$ & $1$ & $-1$ & $3$ & $-3$ & $7$ & $-7$ & $9$ & $-9$ & $21$ & $-21$ & $63$ & $-63$\\ \hline
$28+d$ & $29$ & $27$ & $31$ & $25$ & $35$ & $21$ & $37$ & $19$ & $49$ & $7$ & $91$ & $-35$\\
quadrado perfeito? & não & não & não & \textbf{sim} & não & não & não & não & \textbf{sim} & não & não & não
\end{tabular}
\end{center}

\begin{itemize}
\item $d=-3$: $(n-1)^2=25\Rightarrow n-1=\pm5\Rightarrow n=6$ ou $n=-4$.
\item $d=21$: $(n-1)^2=49\Rightarrow n-1=\pm7\Rightarrow n=8$ ou $n=-6$.
\end{itemize}

\textbf{4) Verificação.}
\begin{itemize}
\item $n=8$: $D=21$ e $2n-1=15$, logo o quociente é $15+3=18$.
\item $n=-6$: $D=21$ e $2n-1=-13$, logo o quociente é $-13+3=-10$.
\item $n=6$: $D=-3$ e $2n-1=11$, logo o quociente é $11-21=-10$.
\item $n=-4$: $D=-3$ e $2n-1=-9$, logo o quociente é $-9-21=-30$.
\end{itemize}
Todos são inteiros. Uma busca computacional em $-2000\le n\le 2000$ confirma que esses são os únicos valores.

\textbf{Resposta:} $n\in\{-6,\,-4,\,6,\,8\}$.

ANSWER: $n\in\{-6,-4,6,8\}$%
\newpage%
\section*{Questão 8}

\section*{Solução}

\textbf{Modelagem.} Em coordenadas, tome $A=(0,0)$, $B=(10,0)$, $C=(10,6)$, $D=(0,6)$. "Inscrito" significa um vértice de $EFGH$ em cada lado de $ABCD$; como $E\in\overline{AB}$ e $F\in\overline{BC}$ são vértices consecutivos, $G\in\overline{CD}$ e $H\in\overline{DA}$. Escreva
\[
E=(x,0),\quad F=(10,y),\quad H=(0,h),\qquad 0<x<10,\ 0<y<6,\ 0<h<6 .
\]

Como $EFGH$ é paralelogramo, $G=F+H-E=(10-x,\;y+h)$. Como $G\in\overline{CD}$ (reta $y=6$), temos $y+h=6$, isto é, $h=6-y$ e $G=(10-x,6)$.

Como o ângulo em $E$ é reto, $\overrightarrow{EF}\cdot\overrightarrow{EH}=0$, com $\overrightarrow{EF}=(10-x,\,y)$ e $\overrightarrow{EH}=(-x,\,h)$:
\[
-x(10-x)+y\,h=0 \;\Longrightarrow\; y(6-y)=x(10-x). \tag{1}
\]
Reciprocamente, dados $x,y$ nos intervalos acima satisfazendo (1), o paralelogramo construído tem um ângulo reto, logo é um retângulo inscrito com os vértices nos lados indicados.

\textbf{a)} Para $0<y<6$, $y(6-y)=9-(y-3)^2$ assume todos os valores em $(0,9]$ (e só eles). Logo (1) tem solução $y\in(0,6)$ se, e somente se, $0<x(10-x)\le 9$. Como $x(10-x)>0$ para $0<x<10$, resta
\[
x^2-10x+9\ge 0 \iff (x-1)(x-9)\ge 0 \iff x\le 1 \ \text{ou}\ x\ge 9 .
\]
Nesse caso $y=3\pm\sqrt{(x-1)(x-9)}$, e $\sqrt{(x-1)(x-9)}=\sqrt{9-x(10-x)}<3$, de modo que $y\in(0,6)$.

Portanto
\[
\boxed{\,0<x\le 1\quad\text{ou}\quad 9\le x<10\,}.
\]
(Verificação: para $x=5$ teríamos $y(6-y)=25>9$, impossível; para $x=1$, $y=3$ e $EFGH$ existe.)

\textbf{b)} Os quatro triângulos retângulos dos cantos têm catetos:
\[
AEH:\ x,\,h;\qquad BEF:\ 10-x,\,y;\qquad CFG:\ 6-y,\,x;\qquad DGH:\ 10-x,\,6-h=y .
\]
Com $h=6-y$, a soma de suas áreas é
\[
\tfrac12\big[x(6-y)+(10-x)y+(6-y)x+(10-x)y\big]=x(6-y)+(10-x)y=6x+10y-2xy .
\]
Logo
\[
[EFGH]=60-6x-10y+2xy .
\]
Impondo $[EFGH]=30$:
\[
6x+10y-2xy=30 \iff xy-3x-5y+15=0 \iff (x-5)(y-3)=0 .
\]
\begin{itemize}
\item $x=5$ é impossível, pois $5\in(1,9)$ contradiz o item (a).
\item Logo $y=3$, e (1) fornece $x(10-x)=9$, ou seja, $x=1$ ou $x=9$ (ambos permitidos por (a)). Então $h=6-y=3$.
\end{itemize}

Para $x=1$, $y=3$, $h=3$:
\[
EF=\sqrt{(10-1)^2+3^2}=\sqrt{90}=3\sqrt{10},\qquad EH=\sqrt{1^2+3^2}=\sqrt{10}.
\]
Para $x=9$ os papéis se invertem: $EF=\sqrt{1+9}=\sqrt{10}$ e $EH=\sqrt{81+9}=3\sqrt{10}$ (é o retângulo simétrico do anterior).

Conferência: $3\sqrt{10}\cdot\sqrt{10}=30$. Além disso, para $x=1$: $E=(1,0)$, $F=(10,3)$, $G=(9,6)$, $H=(0,3)$, com $\overrightarrow{EF}\cdot\overrightarrow{FG}=(9,3)\cdot(-1,3)=0$, confirmando o ângulo reto.

Os lados de $EFGH$ medem $\sqrt{10}$ e $3\sqrt{10}$ (nos casos $x=1$ ou $x=9$).

ANSWER: a) $0<x\le 1$ ou $9\le x<10$; b) os lados medem $\sqrt{10}$ e $3\sqrt{10}$.%
\newpage%
\section*{Questão 9}

\section*{Solução}

\textbf{1. Dados da pirâmide.} Seja $V$ o vértice, $A_1,\dots,A_6$ os vértices do hexágono regular de lado $4$ (circunraio $4$) e $O$ seu centro. Seja $\ell$ a aresta lateral. Pela lei dos cossenos no triângulo $VA_1A_2$, com ângulo $\frac{\pi}{6}$ em $V$,
\[
16=2\ell^2\left(1-\cos\tfrac{\pi}{6}\right)\;\Rightarrow\;\ell^2=\frac{16}{2-\sqrt3}=16\,(2+\sqrt3).
\]
Como $\ell>4$, a pirâmide existe. A pirâmide é reta, logo $VO\perp$ base, e $h=VO$ satisfaz $h^2=\ell^2-4^2=16+16\sqrt3$. O apótema da face é $VM=\sqrt{\ell^2-2^2}=\sqrt{28+16\sqrt3}=2\sqrt{7+4\sqrt3}=2(2+\sqrt3)$, pois $(2+\sqrt3)^2=7+4\sqrt3$.

\textbf{2. Escolha da diagonal e da face.} A pirâmide é invariante pela rotação de $60^\circ$ em torno de $VO$, que leva uma diagonal maior em qualquer outra. Podemos então fixar a diagonal $A_1A_4$.

Seja $\pi$ o plano secante e $F$ a face a que $\pi$ é paralelo. Os planos $\pi$ e $F$ são paralelos, e o plano da base os corta segundo retas paralelas. Logo $\pi\cap\text{base}=A_1A_4$ é paralela à reta que contém a aresta de $F$ na base. Os lados do hexágono paralelos a $A_1A_4$ são apenas $A_2A_3$ e $A_5A_6$. As faces $VA_1A_2$, $VA_6A_1$, $VA_3A_4$ e $VA_4A_5$ ficam excluídas. Por exemplo, o plano de $VA_1A_2$ corta a base na reta $A_1A_2$, que não é paralela a $A_1A_4$.

A reflexão $y\mapsto -y$ fixa $A_1A_4$ e troca $A_2A_3$ com $A_5A_6$. As duas faces dão, portanto, seções congruentes. Tomamos $F=VA_2A_3$. Por uma reta passa um único plano paralelo a $F$, então a seção é única e a área está bem definida.

\textbf{3. O plano $\pi$.} Coordenadas: $O=(0,0,0)$, $A_1=(4,0,0)$, $A_2=(2,2\sqrt3,0)$, $A_3=(-2,2\sqrt3,0)$, $A_4=(-4,0,0)$, $A_5=(-2,-2\sqrt3,0)$, $A_6=(2,-2\sqrt3,0)$, $V=(0,0,h)$.

O plano de $F$ passa por $(0,2\sqrt3,0)$, $(1,2\sqrt3,0)$ e $V$, e sua equação é $\frac{y}{2\sqrt3}+\frac{z}{h}=1$. O plano paralelo a ele que contém o eixo $x$ (a reta $A_1A_4$) é
\[
\pi:\quad h\,y+2\sqrt3\,z=0 .
\]

\textbf{4. Vértices da seção.} Seja $f(x,y,z)=hy+2\sqrt3\,z$. Então
\[
f(A_1)=f(A_4)=0,\qquad f(V)=f(A_2)=f(A_3)=2\sqrt3\,h>0,\qquad f(A_5)=f(A_6)=-2\sqrt3\,h<0 .
\]
Como $f$ é afim, ao longo de uma aresta ela é linear. Examinamos cada aresta:
\begin{itemize}
\item $VA_1$, $VA_4$ e as arestas da base que saem de $A_1$ ou $A_4$: o outro extremo tem $f>0$ ou $f<0$, então $\pi$ só as toca em $A_1$ ou $A_4$.
\item $VA_2$, $VA_3$, $A_2A_3$: $f>0$ nos dois extremos, logo não há interseção. $A_5A_6$ tem $f<0$ nos dois extremos, também sem interseção.
\item $VA_5$ e $VA_6$: $f$ vale $+2\sqrt3\,h$ em $V$ e $-2\sqrt3\,h$ no outro extremo, então se anula no ponto médio. Obtemos $N_5$ e $N_6$, pontos médios de $VA_5$ e $VA_6$.
\end{itemize}
A seção é, portanto, o quadrilátero $A_1N_6N_5A_4$, com lados nas faces $VA_6A_1$, $VA_5A_6$, $VA_4A_5$ e na base.

\textbf{5. Forma e dimensões.} $N_5N_6$ é base média do triângulo $VA_5A_6$, então $N_5N_6=\tfrac12 A_5A_6=2$ e $N_5N_6\parallel A_5A_6\parallel A_1A_4$. A seção é um trapézio de bases $A_1A_4=8$ e $N_5N_6=2$. Ele é isósceles, por simetria na reflexão $x\mapsto -x$.

\emph{Altura.} Seja $M$ o ponto médio de $A_2A_3$ e $M'$ o de $A_5A_6$. Então $O$ é ponto médio de $MM'$ e $N$, ponto médio de $N_5N_6$, é ponto médio de $VM'$. O plano $x=0$ contém $V,M,M',O,N$ e é perpendicular a $A_1A_4$ e a $N_5N_6$. No triângulo $VMM'$, $ON$ liga os pontos médios de $MM'$ e $VM'$, logo $ON\parallel VM$ e $ON=\tfrac12 VM$. $ON$ está em $\pi$ e é perpendicular às duas bases, então é a altura do trapézio:
\[
H=\tfrac12 VM=2+\sqrt3 .
\]
Em coordenadas, $N=(0,-\sqrt3,\tfrac h2)$ e $H=\sqrt{3+\tfrac{h^2}{4}}=\sqrt{7+4\sqrt3}=2+\sqrt3$, o que confirma o resultado.

\textbf{6. Área.}
\[
S=\frac{8+2}{2}\,(2+\sqrt3)=5\,(2+\sqrt3)=10+5\sqrt3 .
\]

\textbf{Conferência.} Numericamente, $h\approx 6{,}6116$ e $H\approx 3{,}7321$, e a área é $S\approx 18{,}66$. A ordem de grandeza é compatível com um trapézio de bases $8$ e $2$ e altura próxima de $3{,}7$. Os vértices obtidos pelo sinal de $f$ coincidem com os da construção sintética.

ANSWER: $S = 5\,(2+\sqrt{3}) = 10+5\sqrt{3}$ (unidades de área)%
\newpage%
\section*{Questão 10}

\section*{Solução}

\textbf{1. Simplificação.} Pelas identidades $1+\cos 2x = 2\cos^2 x$ e $1-\cos 2x = 2\operatorname{sen}^2 x$, e como $\sqrt{u^2}=|u|$,
\[
f(x)=\sqrt{2}\,|\cos x|+\sqrt{2}\,|\operatorname{sen} x|-\operatorname{sen} x-\cos x .
\]
Seja $s=\operatorname{sen}x+\cos x=\sqrt2\,\operatorname{sen}\!\left(x+\tfrac{\pi}{4}\right)$. Como $f$ tem período $2\pi$, basta estudar $x\in(0,\tfrac{\pi}{2})$ (1º quadrante) e $x\in(\pi,\tfrac{3\pi}{2})$ (3º quadrante).

\textbf{2. Primeiro quadrante} ($\operatorname{sen}x>0$, $\cos x>0$). Aqui $|\cos x|+|\operatorname{sen} x| = s$, logo
\[
f(x)=(\sqrt2-1)\,s .
\]
Para $x\in(0,\tfrac{\pi}{2})$ temos $x+\tfrac{\pi}{4}\in(\tfrac{\pi}{4},\tfrac{3\pi}{4})$, então $\operatorname{sen}(x+\tfrac{\pi}{4})\in(\tfrac{\sqrt2}{2},1]$ e $s\in(1,\sqrt2]$. Como $f$ é crescente em $s$ (coeficiente positivo),
\[
f(x)\in\big(\sqrt2-1,\;(\sqrt2-1)\sqrt2\,\big]=\big(\sqrt2-1,\;2-\sqrt2\,\big]\approx(0{,}414;\,0{,}586],
\]
que não contém inteiro algum.

\textbf{3. Terceiro quadrante} ($\operatorname{sen}x<0$, $\cos x<0$). Agora $|\cos x|+|\operatorname{sen} x| = -s$, então
\[
f(x)=-\sqrt2\,s-s=-(\sqrt2+1)\,s .
\]
Para $x\in(\pi,\tfrac{3\pi}{2})$ temos $x+\tfrac{\pi}{4}\in(\tfrac{5\pi}{4},\tfrac{7\pi}{4})$, portanto $\operatorname{sen}(x+\tfrac{\pi}{4})\in[-1,-\tfrac{\sqrt2}{2})$ e $s\in[-\sqrt2,-1)$, ou seja, $-s\in(1,\sqrt2\,]$. A função $x\mapsto -s$ é contínua e percorre todo o intervalo $(1,\sqrt2\,]$ (o valor $\sqrt2$ é atingido em $x=\tfrac{5\pi}{4}$). Assim
\[
f(x)\in\big((\sqrt2+1),\;(\sqrt2+1)\sqrt2\,\big]=\big(\sqrt2+1,\;2+\sqrt2\,\big]\approx(2{,}414;\,3{,}414],
\]
e todo valor desse intervalo é assumido (continuidade). O único inteiro nele é $3$.

\textbf{4. Verificação de $k=3$.} Devemos ter $-s=\dfrac{3}{\sqrt2+1}=3(\sqrt2-1)\approx1{,}243\in(1,\sqrt2\,]$, o que é possível. Explicitamente, $\operatorname{sen}(x+\tfrac{\pi}{4})=-\tfrac{3(2-\sqrt2)}{2}\approx-0{,}879\in[-1,-\tfrac{\sqrt2}{2})$ tem solução com $x+\tfrac{\pi}{4}\in(\tfrac{5\pi}{4},\tfrac{7\pi}{4})$, por exemplo $x=\tfrac{3\pi}{4}+\arcsin\tfrac{3(2-\sqrt2)}{2}\approx3{,}43$ rad, que está em $(\pi,\tfrac{3\pi}{2})$; conferindo numericamente, $f(x)=3$.

\textbf{5. Fronteiras.} Se os eixos fossem incluídos nos quadrantes, os valores adicionais seriam $f(0)=f(\tfrac{\pi}{2})=\sqrt2-1\approx0{,}414$ e $f(\pi)=f(\tfrac{3\pi}{2})=\sqrt2+1\approx2{,}414$, nenhum inteiro. A conclusão não muda.

\textbf{Conclusão.} Nos quadrantes I e III, $f$ assume valores apenas em $(\sqrt2-1,\,2-\sqrt2\,]\cup(\sqrt2+1,\,2+\sqrt2\,]$, e o único inteiro $k$ nesse conjunto é $k=3$.

ANSWER: $k=3$ (único inteiro).%
\newpage%
\section*{Apêndice: parecer do avaliador}%
\subsection*{Questão 1: nota 9,2 (aprovada, rodada 1 de 1)}%
\textbf{Resposta da resolução às cegas do avaliador:} Com 2a + 9r = 87, a razão r é ímpar e |r| $\leq$ 9. Há 10 PAs como sequências ordenadas, ou 5 se a PA e a sua reversa forem identificadas.\newline%
\newline%
Crescentes (r > 0):\newline%
{-} r = 1, a$_{1}$ = 39: 39, 40, \ldots{}, 48\newline%
{-} r = 3, a$_{1}$ = 30: 30, 33, \ldots{}, 57\newline%
{-} r = 5, a$_{1}$ = 21: 21, 26, \ldots{}, 66\newline%
{-} r = 7, a$_{1}$ = 12: 12, 19, \ldots{}, 75\newline%
{-} r = 9, a$_{1}$ = 3: 3, 12, \ldots{}, 84\newline%
\newline%
Decrescentes (r < 0): as reversas das anteriores.\newline%
{-} r = -1, a$_{1}$ = 48\newline%
{-} r = -3, a$_{1}$ = 57\newline%
{-} r = -5, a$_{1}$ = 66\newline%
{-} r = -7, a$_{1}$ = 75\newline%
{-} r = -9, a$_{1}$ = 84\newline%
\newline%
Não há PA constante, pois a = 43,5 não é natural. O resultado é o mesmo se 0 for considerado natural.\par %
A resposta final do candidato coincide com a minha referência, e uma enumeração por força bruta confirmou: 2a$_{1}$+9r=87, r ímpar, |r|$\leq$9, com os 10 pares (a$_{1}$,r) = (39,1), (30,3), (21,5), (12,7), (3,9) e as reversas (48,-1), (57,-3), (66,-5), (75,-7), (84,-9). O enunciado no PDF confere ($\mathbb{N}$=\{1,2,3,\ldots{}\}, 10 naturais, soma 435). A resolução é correta, bem organizada e rigorosa. Mostra r$\in$$\mathbb{Z}$ pela diferença de dois naturais, a paridade de r, o descarte da PA constante (a$_{1}$=43,5), o limite r$\leq$9 (r=11 daria a$_{1}$<0), a verificação a$_{1}$+a$_{1}$$_{0}$=87 e a observação sobre 0$\in$$\mathbb{N}$. Trata as decrescentes por inversão da sequência e lista as razões -1,\ldots{},-9 e os primeiros termos 48,\ldots{},84. As falhas são apenas de apresentação: a família final só traz explicitamente r>0, a contagem 5{-}ou{-}10 fica sem uma leitura canônica, e o argumento de r<0 por simetria é correto, mas sucinto. Nenhuma armadilha foi deixada de lado. Seria publicável como gabarito comentado com pequenos ajustes.%
\begin{itemize}%
\item%
{[}minor{]} A fórmula geral da resposta final (ANSWER) só cobre r $\in$ \{1,3,5,7,9\}. As decrescentes aparecem só por 'leitura inversa', com razões e primeiros termos dados na Conclusão. Seria mais limpo escrever uma família única, a\_n = (87-9r)/2 + (n-1)r com r $\in$ \{$\pm$1, $\pm$3, $\pm$5, $\pm$7, $\pm$9\}. Seria bom também listar explicitamente as 5 decrescentes, por exemplo 48,47,\ldots{},39, no ANSWER.%
\item%
{[}minor{]} A resposta final admite as duas contagens ('5 crescentes, ou 10 contando as decrescentes') sem decidir qual é a leitura canônica. A referência faz o mesmo e a ambiguidade é do enunciado, mas um gabarito comentado poderia dizer que, como sequências ordenadas, são 10 PAs.%
\item%
{[}minor{]} O caso r<0 é tratado apenas pela simetria da inversão (Passo 3). O argumento é válido, porque a inversão é uma bijeção que preserva os termos, a soma e a pertinência a $\mathbb{N}$. Mas ele não mostra diretamente que o menor termo é a$_{1}$$_{0}$ = a$_{1}$+9r $\geq$ 1 $\Rightarrow$ r $\geq$ -9. Uma frase justificando que a inversão é bijetiva e preserva $\mathbb{N}$ tornaria o passo mais transparente.%
\item%
{[}minor{]} O comentário sobre 0 $\in$ $\mathbb{N}$ está correto, mas incompleto. Diz só que r=11 daria a$_{1}$<0. Falta observar que a$_{1}$=0 exigiria 9r=87, sem solução inteira, e que r=9 continua valendo (a$_{1}$=3). O resultado final não é afetado.%
\end{itemize}%
\subsection*{Questão 2: nota 9,5 (aprovada, rodada 1 de 1)}%
\textbf{Resposta da resolução às cegas do avaliador:} S = \{ x $\in$ $\mathbb{R}$ : log$_{2}$(($\surd$5 - 1)/2) < x < 1 \} = ( -log$_{2}$ $\varphi$ , 1 ), com $\varphi$ = (1+$\surd$5)/2. Equivalentemente, log$_{2}$($\surd$5 - 1) - 1 < x < 1, ou seja, aproximadamente -0,694 < x < 1. Os dois extremos ficam de fora, porque nos dois a igualdade vale.\par %
A resposta final do candidato, S = (log$_{2}$(($\surd$5-1)/2), 1) = (-log$_{2}$((1+$\surd$5)/2), 1) $\approx$ (-0,694; 1), coincide com a da minha referência e com a varredura numérica que refiz. Essa varredura, de -10 a 10 com passo 10$^{-}$$^{4}$, deu o intervalo contíguo de -0,6942 a 0,9999. Os valores numéricos citados pelo candidato (1,390>1,38 em x=-0,69; 1,387<1,40 em x=-0,70; log$_{2}$6$\approx$2,58<4) conferem. O argumento é rigoroso e todos os passos são equivalências: o módulo externo é eliminado porque o logaritmo é maior que 1, e a desigualdade vira 2\^{}x+2>2\^{}\{|2x|\}. Em seguida, os dois casos x$\geq$0 e x<0 são tratados com t=2\^{}x, e as restrições de t (t$\geq$1 e 0<t<1) são combinadas ao final. A fatoração (t+1)(t$^2$+t-1) é correta, a raiz negativa é descartada por t>0, e se verifica que ($\surd$5-1)/2 está em (0,1). Os extremos x=1 e x=-log$_{2}$$\varphi$ são corretamente excluídos por igualdade, e x=0 é incluído. Todas as armadilhas da referência foram tratadas. As ressalvas são só de apresentação: um passo algébrico da igualdade no extremo esquerdo não foi escrito, e a varredura numérica foi citada como evidência, mas sem peso lógico.%
\begin{itemize}%
\item%
{[}minor{]} A verificação do extremo esquerdo diz que t$^2$+t=1 implica t+2=t\^{}\{{-}2\}, mas não mostra a passagem. Falta escrever que t$^2$(t+2)=t$^3$+2t$^2$=t(t$^2$+t)+t$^2$=t+t$^2$=1. É um passo de uma linha, mas um gabarito comentado deveria incluí{-}lo.%
\item%
{[}minor{]} A conferência numérica aparece como complemento ('varredura em {[}{-}5,5{]}') e não entra na prova. A justificativa das extremidades e da unicidade dos cruzamentos vem corretamente das equivalências encadeadas e da fatoração, então não há lacuna lógica. Seria melhor não mencionar a varredura, ou dizer que ela é só sanidade, já que ela cobre apenas {[}{-}5,5{]}.%
\item%
{[}minor{]} Não há comentário explícito de que, para x muito negativo, o lado esquerdo tende a 1 e |2x| tende a infinito. Esse comentário daria uma sanidade gráfica e justificaria por que o conjunto é limitado à esquerda. A prova algébrica já cobre isso, pois t>($\surd$5-1)/2 exclui x muito negativo.%
\end{itemize}%
\subsection*{Questão 3: nota 9,0 (aprovada, rodada 1 de 1)}%
\textbf{Resposta da resolução às cegas do avaliador:} Lado do octógono regular: ? = $\surd$(4 - 2$\surd$2)/2 = $\surd$((2 - $\surd$2)/2) = $\surd$2$\cdot$sen($\pi$/8) $\approx$ 0,541. Valor de r: r = 1/16. (Raio da circunferência circunscrita R = 1/$\surd$2; as raízes de z$^{8}$ = r$\cdot$i formam um octógono regular; área = 2$\surd$2$\cdot$R$^2$ = $\surd$2.)\par %
A resolução do candidato está correta e coincide com a minha referência nos dois itens, na ordem pedida (lado $\to$ r). O lado vale ? = $\surd$((2-$\surd$2)/2) = $\surd$(4-2$\surd$2)/2 $\approx$ 0,541 e r = 1/16. Conferi numericamente: com r = 1/16, as oito raízes satisfazem |r + i z$^{8}$| $\approx$ 10$^{-}$$^{1}$$^{7}$. A fórmula do cadarço dá área 1,41421356\ldots{} = $\surd$2 e o lado medido é 0,5411961\ldots{}, igual à forma fechada do candidato.\newline%
\newline%
O desenvolvimento é limpo:\newline%
{-} O sinal foi tratado corretamente: -1/i = i, logo z$^{8}$ = i$\cdot$r.\newline%
{-} De Moivre dá as raízes com módulo R = r\^{}\{1/8\} e argumentos $\pi$/16 + k$\pi$/4, igualmente espaçados, o que justifica o octógono regular. Esse era o ponto fraco apontado na referência, e o candidato o mostrou explicitamente.\newline%
{-} A área foi calculada como 8 triângulos de ângulo central $\pi$/4: A = 2$\surd$2$\cdot$R$^2$, logo R$^2$ = 1/2.\newline%
{-} A passagem de R para r foi feita corretamente: r = R$^{8}$ = (1/2)$^{4}$ = 1/16.\newline%
{-} O lado veio pela lei dos cossenos com ângulo central $\pi$/4, sem confundir com $\pi$/8.\newline%
{-} A conferência independente pela fórmula 2(1+$\surd$2)?$^2$ = $\surd$2 está correta.\newline%
\newline%
Os problemas restantes são só de rigor e polimento: não explicita o polígono convexo em ordem de argumento, não declara que as raízes são distintas e não diz que a rotação por $\pi$/16 é irrelevante para área e lado. Nenhum compromete a correção. É praticamente um gabarito comentado, com pequenas lacunas formais.%
\begin{itemize}%
\item%
{[}minor{]} A resolução não diz que o polígono é o octógono convexo, com os vértices ligados em ordem crescente de argumento. Chama{-}o de "octógono regular inscrito", o que implica isso, mas não explicita que ligar as raízes fora de ordem daria um polígono estrelado de outra área. Bastaria uma frase afirmando isso.%
\item%
{[}minor{]} Não menciona que as 8 raízes são distintas. Os argumentos $\pi$/16 + k$\pi$/4, com k = 0,\ldots{},7, ficam em {[}$\pi$/16, 29$\pi$/16{]} e são todos diferentes módulo 2$\pi$. Sem isso não se garante que há 8 vértices distintos.%
\item%
{[}minor{]} Não comenta que a rotação de $\pi$/16 (ou um eventual erro de sinal em z$^{8}$ = $\pm$ri) não altera área nem lado. O sinal foi tratado corretamente (-1/i = i), então não há erro, apenas uma justificativa que faltou.%
\item%
{[}minor{]} A conferência de |z\_k|$^{8}$ = r é circular, porque apenas repete a definição de R. A conferência útil é a da área pela fórmula 2(1+$\surd$2)?$^2$, que está correta e é independente. A expressão do lado como 2R$\cdot$sen($\pi$/8) = $\surd$2$\cdot$sen($\pi$/8) também seria um complemento natural, mas não é necessária.%
\end{itemize}%
\subsection*{Questão 4: nota 9,0 (aprovada, rodada 1 de 1)}%
\textbf{Resposta da resolução às cegas do avaliador:} a = -2. p(x) = (x - 2)(x + 1)$^2$(x$^2$ + 1)(x$^2$ - 6x + 1). Raízes de p: x = -1 (dupla), x = 2, x = i, x = -i, x = 3 + 2$\surd$2 e x = 3 - 2$\surd$2. Conjunto{-}solução: \{-1, 2, i, -i, 3 + 2$\surd$2, 3 - 2$\surd$2\}.\par %
Resposta final correta e idêntica à minha referência. Conferi em Python puro, sem sympy: a={-}2 (p(2)=170+85a, p({-}1)={-}16{-}8a, p'({-}1)=32+16a, os três nulos só em a={-}2). O produto (x+1)$^2$(x{-}2)(x$^2$+1)(x$^2${-}6x+1) reproduz exatamente os coeficientes de p com a={-}2. s({-}1)=16 e s(2)={-}35. Os erros numéricos em $\pm$i e 3$\pm$2$\surd$2 são desprezíveis. As raízes são {-}1 (dupla), 2, $\pm$i e 3$\pm$2$\surd$2. O candidato fatorou q com a raiz dupla {-}1 e exigiu p'({-}1)=0 além de p(2) e p({-}1), tratando a principal armadilha. Fez a divisão p÷q com quociente e resto corretos e a substituição t=x+1/x com o {-}2 em x$^2$+1/x$^2$=t$^2${-}2, resolvendo t=0 e t=6 e incluindo as raízes complexas $\pm$i. Verificou as multiplicidades por s({-}1)$\neq$0 e s(2)$\neq$0 e deu a conclusão com a fatoração completa. Todas as armadilhas da referência foram evitadas. As fraquezas são de rigor e apresentação: não argumenta explicitamente a suficiência e a unicidade de a, omite várias contas intermediárias (p'({-}1), $\Delta$, s({-}1), s(2)), não faz a conferência de Girard e não expande a fatoração de s. Nenhuma afeta a correção. Nota 9, no limite do que se publicaria como gabarito comentado, com um pouco mais de detalhe nesses passos.%
\begin{itemize}%
\item%
{[}minor{]} A resolução não justifica explicitamente a suficiência e a unicidade de a. p(2)=0, p({-}1)=0 e p'({-}1)=0 são condições necessárias. A suficiência só aparece implicitamente, na divisão exata com resto nulo do passo 3. Faltou uma frase dizendo que cada condição, sendo linear em a, força a={-}2 e que a divisão exata confirma a recíproca. O resto genérico (a+2)(5x$^2$+26x+13) não foi mostrado.%
\item%
{[}minor{]} Vários cálculos aparecem só com o resultado final. É o caso de p'({-}1)=32+16a, sem as parcelas, de s({-}1)=16 e s(2)={-}35, sem as contas, e do discriminante $\Delta$=32 com $\surd$32=4$\surd$2 em x$^2${-}6x+1=0. As contas estão corretas, mas um gabarito comentado deveria mostrá{-}las.%
\item%
{[}minor{]} A fatoração s(x)=(x$^2$+1)(x$^2${-}6x+1) e o produto final dos cinco fatores são apenas afirmados ('Multiplicando os fatores, obtém{-}se\ldots{}'). Faltou um passo de conferência, como a expansão (x$^2$+1)(x$^2${-}6x+1)=x$^{4}${-}6x$^3$+2x$^2${-}6x+1. Também não há checagem por Girard (soma 6, produto 2). A verificação numérica citada só cobre 3$\pm$2$\surd$2.%
\item%
{[}minor{]} No passo 3 o quociente é anunciado antes de ser derivado. A 'conferência' lista os produtos parciais (x$^{4}$q, {-}6x$^3$q, 2x$^2$q, {-}6xq, 1q), mas não mostra a soma coeficiente a coeficiente. A divisão está correta, mas a apresentação é de verificação e não de dedução.%
\item%
{[}minor{]} A multiplicidade 1 de $\pm$i e de 3$\pm$2$\surd$2 não é dita explicitamente. Segue de serem quatro raízes distintas de s, e distintas de {-}1 e 2. Também não há menção a que são 7 raízes contadas com multiplicidade, igual ao grau, como conferência. A lista de sete raízes está lá, mas a observação do grau é implícita.%
\end{itemize}%
\subsection*{Questão 5: nota 9,0 (aprovada, rodada 1 de 1)}%
\textbf{Resposta da resolução às cegas do avaliador:} Abscissa de C = 5$\surd$3 - 9 $\approx$ -0,34. Há dois pontos C possíveis, simétricos em relação ao eixo x, com y = $\pm$$\surd$(180$\surd$3 - 306) $\approx$ $\pm$2,40, e ambos têm essa mesma abscissa. A e B são (1, 2) e (1, -2), com AB = 4 e AC + BC = 6.\par %
A resposta final 5$\surd$3 - 9 $\approx$ -0,34 está correta e coincide com a minha referência. Conferi numericamente: com x = 5$\surd$3 - 9 e y = $\pm$2,40190, AC $\approx$ 1,39873, BC $\approx$ 4,60127 e o perímetro dá 10,0000. Uma varredura de 400 000 pontos sobre a elipse dada acha apenas os dois pontos (-0,3397; $\pm$2,4019), e o perímetro varia de 8 a $\approx$ 10,77, então o valor 10 é de fato atingido. Reli o enunciado no PDF (Questão 05) e ele confere com a transcrição. A resolução é limpa e bem estruturada. Encontra A=(1,2) e B=(1,-2) e reduz o perímetro a AC+BC=6. Em seguida identifica a elipse auxiliar (x-1)$^2$/5 + y$^2$/9 = 1 (centro (1,0), a'=3, c'=2, b'$^2$=5, eixo vertical) e a menciona com a condição de existência 6 > 4. A interseção dá x$^2$+18x+6=0, com raízes -9 $\pm$ 5$\surd$3. A raiz -9-5$\surd$3 é descartada por |x| $\leq$ $\surd$3, e o candidato nota que y$^2$>0 e que x$\neq$1 (triângulo não degenerado). As contas conferem (x$^2$ = 156 - 90$\surd$3, 9(x-1)$^2$+5(6-2x$^2$)=45). Das armadilhas que listei, o candidato trata praticamente todas. A de reversibilidade não se aplica a este método, pois o lugar geométrico é uma equivalência. Os defeitos são só de acabamento: verificação exata ausente, y exato não dado e WLOG de A/B pouco justificado. Não há erro matemático. Está no nível de um gabarito comentado, com pequenas melhorias possíveis.%
\begin{itemize}%
\item%
{[}minor{]} A verificação do passo 6 é só numérica, com valores arredondados. Faltou uma checagem exata (por exemplo, y$^2$ = 180$\surd$3 - 306, ou AC = 3 - 2y/3 e BC = 3 + 2y/3). O argumento já está completo sem ela, porque o lugar geométrico PA+PB=6 é equivalência e não só implicação. Mas o candidato não diz isso, e o leitor pode se perguntar se elevar ao quadrado criou raízes espúrias.%
\item%
{[}minor{]} O candidato não dá o valor exato de y ($\pm$$\surd$(180$\surd$3 - 306)), só y$^2$ $\approx$ 5,77. Também não verifica que C $\neq$ A e C $\neq$ B, embora isso decorra de x\_C $\neq$ 1.%
\item%
{[}minor{]} O WLOG de A=(1,2) e B=(1,-2) é dito com 'podemos tomar' e sem justificativa explícita. Bastava uma frase: os rótulos são intercambiáveis, e a troca equivale à simetria em relação ao eixo x, que mantém a abscissa de C.%
\item%
{[}minor{]} A resposta não comenta que a abscissa é única apesar de haver dois pontos C ($\pm$y). O candidato diz que há dois pontos com essa abscissa, o que é suficiente. Também não nota que C fica do lado oposto a A e B em relação ao eixo y, observação opcional.%
\end{itemize}%
\subsection*{Questão 6: nota 9,2 (aprovada, rodada 1 de 1)}%
\textbf{Resposta da resolução às cegas do avaliador:} 235 sêxtuplas. Pela convenção da prova, $\mathbb{N}$ = \{1, 2, 3, \ldots{}\}. Todas têm x$_{2}$ = 6 e x$_{1}$ = 8 - 2x$_{6}$, com x$_{6}$ $\in$ \{1, 2, 3\} e x$_{3}$ + x$_{4}$ + x$_{5}$ = 16 - x$_{6}$. Contagem: 91 + 78 + 66 = 235.\par %
Resposta final 235, que coincide com a referência e com a minha conferência. Refiz por força bruta (python3, rodado na pasta de trabalho): com $\mathbb{N}$=\{1,2,3,\ldots{}\} há 235 sêxtuplas, distribuídas em (x6=1, x1=6, x2=6): 91; (x6=2, x1=4, x2=6): 78; (x6=3, x1=2, x2=6): 66. Com $\mathbb{N}$$_{0}$ seriam 605, valor incorreto para este enunciado. O enunciado original (Questão 06) bate com a transcrição e a convenção $\mathbb{N}$=\{1,2,3,\ldots{}\} vem do cabeçalho. A resolução está correta e bem organizada. Adota a convenção certa, agrupa x3+x4+x5=S, resolve o sistema com a álgebra correta (x2=6, x1=8-2x6, S=16-x6) e confere a conta de (III), que dá 60+S+x6. Limita x6 a \{1,2,3\}, aplica corretamente C(S-1,2) e soma 91+78+66=235 sem erro aritmético. Os defeitos são só de rigor e de comentário. Falta declarar que a redução é uma equivalência, justificar estrelas e barras, citar explicitamente o caso x6=4 (x1=0) e a variante $\mathbb{N}$$_{0}$ (605), e dar uma verificação mais completa, com uma sêxtupla explícita e a distribuição da força bruta. Nada disso compromete a correção. É um gabarito sólido, quase publicável, ao qual faltam esses pequenos acabamentos.%
\begin{itemize}%
\item%
{[}minor{]} A equivalência entre o sistema original e o reduzido não é declarada. Os passos (II)-(I), a substituição em (I) e a substituição em (III) são todos reversíveis, então a redução é de fato equivalente. O texto só mostra a direção 'sistema $\Rightarrow$ valores' e não diz que a recíproca vale. Isso importa porque a contagem exige que toda terna com soma S forneça uma solução do original. Falta uma frase como 'as passagens são reversíveis, logo toda (x1,x2,x3,x4,x5,x6) com x2=6, x1=8-2x6 e x3+x4+x5=16-x6 resolve o sistema original'.%
\item%
{[}minor{]} A armadilha do zero e do caso x6=4 não é tratada de forma explícita. O texto afirma x6 $\leq$ 3 porque x1 = 8-2x6 $\geq$ 1, o que está correto, mas não diz que x6=4 daria x1=0, excluído pela convenção $\mathbb{N}$=\{1,2,\ldots{}\}. Também não comenta que a variante com $\mathbb{N}$$_{0}$ daria 605 (conferi por força bruta: 605). Um gabarito comentado destacaria isso, já que é o erro mais provável.%
\item%
{[}minor{]} Estrelas e barras é citado sem justificativa. Bastaria a substituição yi = xi - 1 $\geq$ 0, que leva a y3+y4+y5 = S-3 e ao número C(S-3+2, 2) = C(S-1, 2). A fórmula está correta e os valores 91, 78 e 66 estão certos, mas a passagem é só afirmada.%
\item%
{[}minor{]} A verificação numérica está incompleta. O texto verifica só x6=1, x1=x2=6 e S=15, sem exibir uma sêxtupla concreta, como (6,6,5,5,5,1). A 'contagem por força bruta' é só mencionada, sem o código ou a distribuição por x6. Seria mais convincente exibir também a distribuição (91, 78, 66) obtida no computador ou checar um caso de borda, como x6=3 com x1=2.%
\item%
{[}minor{]} O texto diz que as 3 equações em 4 incógnitas (x1, x2, S, x6) deixam 'um grau de liberdade'. Isso vale para o sistema em S, mas o sistema original tem 3 graus de liberdade nas 6 variáveis: x6 e a escolha da terna com soma S. A frase, como está, pode confundir o leitor. Convém esclarecer que a multiplicidade vem da composição de S.%
\end{itemize}%
\subsection*{Questão 7: nota 9,3 (aprovada, rodada 1 de 1)}%
\textbf{Resposta da resolução às cegas do avaliador:} n $\in$ \{-6, -4, 6, 8\}. Os quocientes valem, respectivamente, -10, -30, -10 e 18. A chave é Q(n) = 2n - 1 + 63/(n$^2$ - 2n - 27), então é preciso que n$^2$ - 2n - 27 divida 63.\par %
Resposta final n $\in$ \{-6, -4, 6, 8\}, igual à da minha referência. Conferi por força bruta em |n| $\leq$ 5000: o conjunto é exatamente \{-6, -4, 6, 8\}, com quocientes -10, -30, -10 e 18, os mesmos valores que o candidato obtém. O enunciado transcrito confere com o PDF (Questão 07, p. 8 do arquivo).\newline%
\newline%
A resolução está correta e quase de nível de gabarito comentado:\newline%
{-} Mostra que D(n) = (n-1)$^2$ - 28 nunca se anula para n inteiro, pois 28 não é quadrado perfeito.\newline%
{-} Faz a divisão polinomial sem erro: N = (2n-1)D + 63. Eu refiz a conta e o resto é 63.\newline%
{-} Reduz a D | 63 apenas pela identidade N = (2n-1)D + 63, sem usar Euclides com passos falsos.\newline%
{-} Testa os 12 divisores, positivos e negativos, via (n-1)$^2$ = 28 + d. Acha d = -3 (n = 6 ou -4) e d = 21 (n = 8 ou -6), com os dois sinais de n-1.\newline%
{-} Verifica os quatro valores diretamente.\newline%
\newline%
Todas as armadilhas listadas foram tratadas: denominador não nulo, resto 63, divisores negativos, dois sinais de cada raiz, argumento de finitude via divisores de 63 e verificação final. Os problemas restantes são só de apresentação: a citação da busca computacional como 'confirmação', e a justificativa um pouco enxuta da equivalência Q inteiro $\Leftrightarrow$ D | 63 e do descarte de d = -63. Nenhum afeta a correção. Nota 9.3.%
\begin{itemize}%
\item%
{[}minor{]} O fechamento cita uma busca computacional em {-}2000 $\leq$ n $\leq$ 2000 como confirmação de que não há outras soluções. Isso não pertence a uma resolução discursiva e não prova nada. A unicidade já decorre, de forma rigorosa, de D(n) dividir 63 (12 valores possíveis) e de o teste de quadrado perfeito ter sido feito para todos os 12. Recomendo remover a frase ou trocá{-}la por uma frase explícita de finitude, por exemplo: 'como |D| $\leq$ 63, vem (n-1)$^2$ $\leq$ 91, e os casos foram esgotados'.%
\item%
{[}minor{]} A equivalência 'Q inteiro $\Leftrightarrow$ D | 63' é enunciada de forma muito curta. Faltou dizer que 63/D = Q - (2n-1) é diferença de inteiros no sentido $\to$, e que no sentido ? 63/D inteiro implica Q inteiro. Faltou também explicitar que 'divide' vale em $\mathbb{Z}$, com divisores negativos. A lista $\pm$1,\ldots{},$\pm$63 mostra que o candidato sabe disso, então o prejuízo é só de clareza.%
\item%
{[}minor{]} O candidato diz que (n-1)$^2$ = 28 + d deve ser 'quadrado perfeito (e não negativo)', mas não diz explicitamente que d = -63 dá -35 < 0 e por isso é descartado. A tabela mostra -35 com 'não', então a conclusão está correta, mas a justificativa fica implícita. Também não comenta que 28 + d é um quadrado perfeito $\Rightarrow$ n inteiro, ou seja, que as raízes n = 1 $\pm$ k são inteiras. É trivial, mas é o passo que fecha a equivalência.%
\end{itemize}%
\subsection*{Questão 8: nota 9,3 (aprovada, rodada 1 de 1)}%
\textbf{Resposta da resolução às cegas do avaliador:} a) 0 < x $\leq$ 1 ou 9 $\leq$ x < 10, ou seja, x $\in$ (0,1{]} ? {[}9,10). A condição é x(10-x) = y(6-y), com y = BF.\newline%
b) Os lados medem $\surd$10 e 3$\surd$10. Isso ocorre com BF = 3 e x = 1 ou x = 9, que são configurações espelhadas com os mesmos lados.\par %
As respostas finais do candidato coincidem com as da minha referência e estão corretas. (a) 0<x$\leq$1 ou 9$\leq$x<10. (b) lados $\surd$10 e 3$\surd$10, com y=3 e x=1 ou 9 (retângulos espelhados). Conferi o enunciado no PDF original, que bate com a transcrição. Refiz a conta em Python puro (sem numpy): para x em (0,1{]}?{[}9,10), y=3$\pm$$\surd$((x-1)(x-9)) dá um retângulo inscrito com ângulos retos, y$\in$(0,6) e G=(10-x,6), H=(0,6-y). A fórmula S=30+2(x-5)(y-3) vale em todos os casos testados. A área 30 só aparece em x=1 e x=9, com lados 3$\surd$10 e $\surd$10. A modelagem do candidato (E=(x,0), F=(10,y), H=(0,h), G=F+H-E, h=6-y) é limpa e declara a hipótese de inscrição (G$\in$CD, H$\in$DA). A condição de ângulo reto dá y(6-y)=x(10-x). A volta (suficiência) está feita: y(6-y) cobre (0,9{]}, e y=3$\pm$$\surd$(9-x(10-x)) fica estritamente em (0,6), então os extremos x=1 e x=9 entram e x=0 e x=10 são excluídos pelo enunciado. A parte (b) usa a soma dos quatro triângulos de canto (catetos e áreas conferem) e fatora (x-5)(y-3)=0 em vez de dividir. Descarta x=5 pelo item (a), calcula os lados com os dois valores de x e confere com 3$\surd$10$\cdot$$\surd$10=30 e com o produto escalar. Restam só lacunas menores de rigor e de acabamento.%
\begin{itemize}%
\item%
{[}minor{]} A atribuição G$\in$CD e H$\in$DA é declarada, mas justificada só por 'E e F são consecutivos'. Não se diz por que a outra atribuição (G em DA e H em CD) não dá um retângulo convexo inscrito, e o texto não menciona que a leitura 'um vértice por lado' é a interpretação de 'inscrito'. A leitura padrão é aceita, então é uma lacuna pequena.%
\item%
{[}minor{]} Não se comenta que, para x em (0,1)?(9,10), existem dois retângulos (y e 6-y), e que x=1 ou 9 dá um só (y=3). Também não se mostra explicitamente que H=(0,6-y) e G=(10-x,6) ficam no interior dos lados. Isso é imediato pelas faixas dadas, mas um gabarito comentado diria uma frase a respeito. Também não é dito que a escolha da raiz $\pm$ não altera o conjunto de x possíveis.%
\item%
{[}minor{]} O candidato não mostra a interpretação geométrica (semelhança AEH\textasciitilde{}BFE) nem a observação de que o centro de EFGH coincide com o de ABCD. Não é exigido, e o passo analítico basta. Faltou também registrar que a fórmula da área S=30+2(x-5)(y-3) varia, de modo que a área 30 não vale para todo retângulo inscrito.%
\item%
{[}minor{]} A resposta de (b) poderia trazer os valores aproximados ($\approx$3,16 e $\approx$9,49) e dizer explicitamente qual lado é EF e qual é FG em cada caso. O candidato faz isso para x=1 e x=9, mas a linha final 'lados medem $\surd$10 e 3$\surd$10' é um pouco seca.%
\end{itemize}%
\subsection*{Questão 9: nota 9,4 (aprovada, rodada 2 de 2)}%
\textbf{Resposta da resolução às cegas do avaliador:} Área da seção = 5(2 + $\surd$3) = 10 + 5$\surd$3 $\approx$ 18,66. A seção é um trapézio isósceles de bases 8 (a diagonal maior) e 2 (base média da face oposta, com vértices nos pontos médios de duas arestas laterais) e altura 2 + $\surd$3.\par %
A resposta final do candidato, 5(2+$\surd$3) = 10+5$\surd$3 $\approx$ 18,66, coincide com a minha referência. Conferi numericamente em Python puro: plano por A0A3 paralelo a VA1A2, interseção com as 12 arestas, 4 vértices, área 18,6603 = 10+5$\surd$3, h $\approx$ 6,6116 e H = $\surd$(3+h$^2$/4) $\approx$ 3,7321 = 2+$\surd$3. O enunciado transcrito confere com o original (Q09, ITA 2027), e o ângulo de $\pi$/6 no vértice foi verificado (30$^\circ$). Todos os passos do candidato são válidos: ?$^2$ = 16(2+$\surd$3) > 16 (a pirâmide existe), h$^2$ = 16(1+$\surd$3), VM = 2(2+$\surd$3) por $\surd$(7+4$\surd$3) = 2+$\surd$3. A escolha da diagonal por simetria está correta, e só A2A3 e A5A6 são paralelas a A1A4, o que exclui as outras quatro faces. O plano é h$\cdot$y+2$\surd$3$\cdot$z = 0. Os sinais de f (zero em A1, A4; positivo em V, A2, A3; negativo em A5, A6) mostram que VA2 e VA3 não são cortadas e que VA5 e VA6 são cortadas nos pontos médios. A seção é o trapézio isósceles A1N6N5A4 de bases 8 e 2. A altura vem de ON = VM/2 com ON ? bases, e é confirmada por coordenadas. A área é (8+2)/2$\cdot$(2+$\surd$3). Todas as armadilhas da referência foram tratadas: faces válidas e simetria, fator 1/2 da base média, altura $\neq$ perna $\neq$ h/2, desaninhar o radical, forma exata e justificativa. Os problemas restantes são só de apresentação: o uso de $\pi$ para o plano, um passo da lei dos cossenos omitido, a convexidade pressuposta e uma 'conferência' um pouco solta. É praticamente um gabarito comentado publicável.%
\begin{itemize}%
\item%
{[}minor{]} O plano secante é chamado de $\pi$, símbolo que também aparece como ângulo $\pi$/6 no enunciado. Isso pode confundir o leitor. Seria melhor usar $\alpha$ ou $\Pi$ para o plano.%
\item%
{[}minor{]} A passagem 'interseção do plano com o poliedro convexo = polígono de vértices nas interseções com as arestas' é usada sem ser enunciada. A lista de arestas por sinal de f cobre todas as 12 arestas e está correta, mas uma frase de justificativa (convexidade, f afim) daria mais rigor.%
\item%
{[}minor{]} Na lei dos cossenos, o passo 1-cos($\pi$/6) = (2-$\surd$3)/2 é omitido, e VM = 2/tg($\pi$/12) não é apresentado como verificação cruzada. As contas estão certas: ?$^2$ = 16(2+$\surd$3), h$^2$ = 16(1+$\surd$3), VM = 2(2+$\surd$3). Faltou só mostrar o passo intermediário.%
\item%
{[}minor{]} A 'Conferência' diz que os vértices por sinal de f coincidem com 'a construção sintética'. Essa construção sintética completa (plano paralelo a VM por O, base média em VA5A6) não aparece antes. O argumento sintético só entra no cálculo da altura. É uma afirmação de verificação um pouco solta.%
\item%
{[}minor{]} A exclusão das faces é correta e completa, mas poderia dizer explicitamente que as duas arestas da base paralelas a A1A4 são A2A3 e A5A6. Poderia dizer também que, para as demais faces, o plano por A1 paralelo à face coincidiria com a própria face, que não contém A4. O candidato usa só o argumento das retas na base, que basta.%
\end{itemize}%
\subsection*{Questão 10: nota 9,2 (aprovada, rodada 1 de 1)}%
\textbf{Resposta da resolução às cegas do avaliador:} k = 3 (único inteiro). No 1º quadrante f $\in$ ($\surd$2-1, 2-$\surd$2{]} $\approx$ (0,41; 0,59{]}, sem inteiros; no 3º quadrante f $\in$ ($\surd$2+1, 2+$\surd$2{]} $\approx$ (2,41; 3,41{]}, que contém apenas o inteiro 3, atingido em duas soluções (x $\approx$ 1,09$\pi$ e x $\approx$ 1,41$\pi$).\par %
A resposta final k = 3 (único inteiro) está correta e coincide com a minha referência. Conferi numericamente: em {[}0, $\pi$/2{]} a imagem de f é {[}0,4142; 0,5858{]}, em {[}$\pi$, 3$\pi$/2{]} é {[}2,4142; 3,4142{]}, e x = 3$\pi$/4 + arcsin(3(2-$\surd$2)/2) $\approx$ 3,4293 está em ($\pi$, 3$\pi$/2) com f(x) = 3,000000. A resolução trata corretamente a armadilha principal ($\surd$(cos$^2$x) = |cos x|), obtendo f = ($\surd$2-1)s no 1º quadrante e f = -($\surd$2+1)s no 3º. As imagens de s = $\surd$2$\cdot$sen(x+$\pi$/4) estão certas: (1, $\surd$2{]} no 1º e {[}-$\surd$2, -1) no 3º. Ela justifica por continuidade que todo valor do intervalo é atingido, verifica explicitamente que k = 3 é atingido e comenta as fronteiras (eixos: $\surd$2-1 e $\surd$2+1, não inteiros). Não há erro de conta nem de lógica. As falhas são só de acabamento: os limites do intervalo são justificados por decimais em vez de desigualdades exatas, e a segunda solução simétrica não é mencionada. A resolução é clara, completa e próxima de um gabarito comentado.%
\begin{itemize}%
\item%
{[}minor{]} A ausência de inteiros na imagem do 1º quadrante e a unicidade do 3 no 3º quadrante se apoiam só em aproximações decimais ($\approx$0,414--0,586 e $\approx$2,414--3,414). Faltou a justificativa exata: 0 < $\surd$2-1 e 2-$\surd$2 < 1, e 2 < $\surd$2+1 e 2+$\surd$2 < 4. Bastava observar que 1 < $\surd$2 < 2. As desigualdades são triviais, mas um gabarito comentado deveria escrevê{-}las.%
\item%
{[}minor{]} Na verificação de k=3 aparece apenas um valor de x ($\approx$3,43), obtido por 3$\pi$/4 + arcsin(...). Não é dito que x+$\pi$/4 = $\pi$ + arcsin(c) cai em (5$\pi$/4, 7$\pi$/4), nem que há uma segunda solução simétrica em torno de 5$\pi$/4 (x$\approx$4,42). A existência é garantida pelo argumento de continuidade, então a omissão não afeta a correção. A conferência numérica foi só citada ("conferindo numericamente, f(x)=3"), sem mostrar a conta.%
\item%
{[}minor{]} Não comenta que, nos 2º e 4º quadrantes, f assume valores como 1 e 2, que ficam fora do pedido. Também não explica por que a identidade \textbackslash{}|$\cdot$\textbackslash{}| muda a expressão de f de um quadrante para outro. Os pontos centrais, o uso do módulo e a mudança de sinal de f nos dois quadrantes, estão corretos e explícitos. Seria apenas um comentário didático adicional.%
\end{itemize}%
\end{document}